\chapter{Connecting the Android Application}

\section{Install the Application}

Install the RC Controller Android application on your smartphone.

\section{Enable Bluetooth}

Make sure Bluetooth is enabled on your phone.

\section{Connect to the ESP32}

Open the RC Controller App and select the ESP32 device from the list of available Bluetooth devices.

Once connected, the control interface will appear on the screen.

\section{Connection Status}

The application will display the connection status.

If the connection is successful, control commands can be sent to the ESP32.

\section{Debug Mode (Viewing Bluetooth Communication)}

The firmware includes an optional debug system that allows users to monitor the communication between the Android application and the ESP32.

This feature is useful when troubleshooting connection problems or verifying that commands from the application are being received correctly.

\subsection{Enable Debug Mode}

Open the file:

\begin{verbatim}
	debug_config.h
\end{verbatim}

Inside this file you will find a configuration flag for enabling debug output.

Enable debugging by setting the flag to:

\begin{verbatim}
	#define DEBUG_ENABLED 1
\end{verbatim}

If debugging is disabled, the value will typically be:

\begin{verbatim}
	#define DEBUG_ENABLED 0
\end{verbatim}

After changing this value, upload the firmware again to the ESP32.

\subsection{Open the Serial Monitor}

Open the Arduino Serial Monitor:

\begin{verbatim}
	Tools -> Serial Monitor
\end{verbatim}

Set the baud rate to:

\begin{verbatim}
	115200
\end{verbatim}

Restart the ESP32 to view the startup messages.

\subsection{Debug Output}

When debug mode is enabled the ESP32 will print diagnostic messages to the Serial Monitor.

Typical messages may include:

\begin{verbatim}
	[DEBUG] Bluetooth initialized
	[DEBUG] Waiting for client connection
	[DEBUG] Client connected
\end{verbatim}

When commands are received from the Android application, the ESP32 may display messages similar to:

\begin{verbatim}
	[DEBUG] Received command: BTN_A
	[DEBUG] Joystick X: 120
	[DEBUG] Joystick Y: 85
\end{verbatim}

These messages confirm that the Bluetooth communication between the smartphone and the ESP32 is working correctly.

\subsection{Troubleshooting Using Debug Mode}

If the Android application connects but the robot does not respond:

\begin{itemize}
	\item Verify that debug messages appear when pressing buttons in the application.
	\item If no messages appear, the ESP32 may not be receiving Bluetooth data.
	\item Check that the correct ESP32 device is selected in the application.
	\item Verify that the firmware was uploaded successfully.
\end{itemize}

If debug messages appear but the hardware does not respond, the issue is likely related to the wiring or pin configuration.

\section{Using Debug Mode to Verify Communication}

The firmware includes a built-in debugging system that allows the user to see the data being received from the Android application.

This is extremely useful when verifying that the Bluetooth connection is working correctly and that the ESP32 is receiving control commands.

All debug messages are displayed in the Arduino IDE Serial Monitor.

\subsection{Step 1: Enable Debug Mode}

Open the file:

\begin{verbatim}
	debug_config.h
\end{verbatim}

Locate the following line:

\begin{verbatim}
	#define DEBUG_ENABLED 1
\end{verbatim}

Set the value to:

\begin{itemize}
	\item \texttt{1} to enable debug output
	\item \texttt{0} to disable debug output
\end{itemize}

When debugging is enabled, the firmware will print control data received from the Android application.

Upload the firmware again after making any changes.

\subsection{Step 2: Open the Serial Monitor}

In the Arduino IDE open:

\begin{verbatim}
	Tools -> Serial Monitor
\end{verbatim}

Set the baud rate to:

\begin{verbatim}
	115200
\end{verbatim}

Restart the ESP32 or press the reset button.

The Serial Monitor will begin displaying debug messages.

\subsection{Step 3: Move Controls in the App}

After connecting the Android application to the ESP32, move the controls on the screen:

\begin{itemize}
	\item Move the joysticks
	\item Rotate knobs
	\item Toggle switches
	\item Press buttons
\end{itemize}

The ESP32 will print the received data in the Serial Monitor.

\section{Debug Output Types}

The debug system can print different categories of information depending on the configuration in \texttt{debug\_config.h}.

\subsection{Joystick Debug (Sticks)}

If the following option is enabled:

\begin{verbatim}
	#define DBG_STICKS 1
\end{verbatim}

The ESP32 will print the joystick positions when they change.

Example output:

\begin{verbatim}
	L Stick X: 512    L Stick Y: 510
	L Stick X: 650    L Stick Y: 480
	
	R Stick X: 300    R Stick Y: 700
\end{verbatim}

These values represent the position of the left and right joystick on the Android application.

Typical range:

\begin{itemize}
	\item Minimum: around 0
	\item Center: around 512
	\item Maximum: around 1023
\end{itemize}

If the numbers change when you move the joystick, the communication is working correctly.

\subsection{Knob Debug}

If the following option is enabled:

\begin{verbatim}
	#define DBG_KNOBS 1
\end{verbatim}

The ESP32 will print knob values.

Example output:

\begin{verbatim}
	Left Knob: 340
	Left Knob: 512
	Left Knob: 890
	
	Right Knob: 120
	Right Knob: 780
\end{verbatim}

These values correspond to the knob rotation in the Android application.

Knobs can be used to control parameters such as:

\begin{itemize}
	\item Motor speed
	\item Servo angle
	\item Light intensity
	\item Camera tilt
\end{itemize}

\subsection{Switch Debug}

If the following option is enabled:

\begin{verbatim}
	#define DBG_SWITCHES 1
\end{verbatim}

The ESP32 will display the switch status.

Example output:

\begin{verbatim}
	Switch Byte: 0x05
	S1 = ON
	S2 = OFF
	S3 = ON
	S4 = OFF
	S5 = OFF
	S6 = OFF
\end{verbatim}

Each switch corresponds to a toggle button in the Android application.

These switches are commonly used for:

\begin{itemize}
	\item Turning lights on and off
	\item Activating relays
	\item Changing driving modes
	\item Enabling autonomous functions
\end{itemize}

\subsection{Event Debug}

If the following option is enabled:

\begin{verbatim}
	#define DBG_EVENTS 1
\end{verbatim}

The ESP32 will print a message when an event packet arrives.

Example output:

\begin{verbatim}
	--- EVENT PACKET (PRESS) ---
	Event ID: 0x03
	Checksum: 0x7A
	--------------
	
\end{verbatim}

Events are triggered when the user presses a button in the Android application.

Typical uses include:

\begin{itemize}
	\item Activating a special function
	\item Resetting a system
	\item Triggering a sound or animation
	\item Launching an action on the robot
\end{itemize}

\section{Debugging the Telemetry System}

The firmware includes a telemetry system that allows the ESP32 to send
real-time information back to the Android application.

Telemetry can include:

\begin{itemize}
	\item Numeric values
	\item Indicators
	\item Real-time plots
\end{itemize}

For development and testing, the firmware provides a telemetry debug
system that allows the user to simulate telemetry values using the
Arduino Serial Monitor.

This system is configured in the file:

\begin{center}
	\texttt{debug\_config.h}
\end{center}

Inside this file you can select the telemetry debug mode.

\begin{verbatim}
	#define DBG_NONE        0
	#define DBG_PANEL       1
	#define DBG_INDICATOR   2
	#define DBG_PLOT        3
	
	#define TELEMETRY_DEBUG_MODE DBG_INDICATOR
\end{verbatim}

Only one debug mode should be active at a time.

\subsection{Normal Operation}

\begin{verbatim}
	#define TELEMETRY_DEBUG_MODE DBG_NONE
\end{verbatim}

In this mode the ESP32 reads telemetry values from real hardware sensors.

Examples include:

\begin{itemize}
	\item analog sensors
	\item battery voltage measurement
	\item digital inputs
\end{itemize}

This mode should be used when the robot is fully built.

\subsection{Numeric Panel Debug Mode}

\begin{verbatim}
	#define TELEMETRY_DEBUG_MODE DBG_PANEL
\end{verbatim}

This mode allows testing the numeric telemetry panels without connecting
any sensors.

When this mode is enabled, values entered in the Arduino Serial Monitor
are parsed and sent to the Android application as numeric panel telemetry.

\subsubsection{Serial Monitor Input Format}

Open the Arduino Serial Monitor and send values using the following format:

\begin{verbatim}
	LEFT_VALUE,RIGHT_VALUE
\end{verbatim}

Both values may include a decimal part.

Example inputs:

\begin{verbatim}
	12.3,45.8
	120.0,350.5
	5,10
\end{verbatim}

\subsubsection{Value Conversion}

The values entered in the Serial Monitor are internally converted to the
telemetry format used by the firmware.

Each value is multiplied by 10 and converted to an integer before being
transmitted to the Android application.

Examples:

\begin{itemize}
	\item 12.3 becomes 123
	\item 45.8 becomes 458
	\item 120.0 becomes 1200
\end{itemize}

On the Android side, the received integer value is divided by 10 to
display the correct decimal value on the numeric panel.

\subsubsection{Valid Range}

The allowed range for each numeric value is:

\begin{itemize}
	\item 0.0 to 6553.5
\end{itemize}

This range corresponds to the maximum value that can be transmitted
using a 16-bit unsigned integer with one decimal place of precision.

\subsubsection{Example}

If the following value is entered in the Serial Monitor:

\begin{verbatim}
	123.4,56.7
\end{verbatim}

The firmware will transmit:

\begin{verbatim}
	1234,567
\end{verbatim}

The Android application will then display:

\begin{verbatim}
	123.4   56.7
\end{verbatim}

on the left and right numeric panels respectively.

This debug mode is useful for verifying the correct operation of the
telemetry transmission and the Android user interface before connecting
real sensors.


\subsection{Indicator Debug Mode}

\begin{verbatim}
	#define TELEMETRY_DEBUG_MODE DBG_INDICATOR
\end{verbatim}

This mode allows testing indicator widgets.

Send values through the Serial Monitor using this format:

\begin{verbatim}
	ANALOG_VALUE,BATTERY_LEVEL,DIGITAL_MASK
\end{verbatim}

Example:

\begin{verbatim}
	45,80,5
\end{verbatim}

Value ranges:

\begin{itemize}
	\item Analog value: 0 to 100
	\item Battery level: 0 to 100
	\item Digital mask: 0 to 255
\end{itemize}

The digital mask represents the state of digital indicators using bit flags.

Example values:

\begin{verbatim}
	1
	2
	4
	8
\end{verbatim}

Each value activates a different indicator.

Multiple indicators can be activated by combining values.

Example:

\begin{verbatim}
	3
\end{verbatim}

Binary representation:

\begin{verbatim}
	00000011
\end{verbatim}

This activates the first two indicators.

\subsection{Plot Debug Mode}

\begin{verbatim}
	#define TELEMETRY_DEBUG_MODE DBG_PLOT
\end{verbatim}

This mode simulates three telemetry signals that are displayed as
real-time graphs in the Android application.

Send values using the Serial Monitor:

\begin{verbatim}
	VALUE1,VALUE2,VALUE3
\end{verbatim}

Example:

\begin{verbatim}
	10,40,80
\end{verbatim}

Valid range:

\begin{itemize}
	\item 0 to 255
\end{itemize}

These values will appear in the telemetry plots of the application.

\subsection{How the Telemetry Debug System Works}

The debug values are read by the module:

\begin{center}
	\texttt{telemetry\_source.cpp}
\end{center}

This module replaces real sensor values with values received through the
Serial interface when a debug mode is active.

Telemetry packets are created and transmitted by:

\begin{center}
	\texttt{telemetry.cpp}
\end{center}

This module builds telemetry packets and sends them to the Android
application using Bluetooth.

This architecture allows the telemetry interface to be tested without
connecting any physical sensors.


\section{Why Debug Mode is Important}

Debug mode allows users to quickly verify that:

\begin{itemize}
	\item The Android application is connected
	\item Bluetooth data is being received
	\item Control values are changing correctly
\end{itemize}

If values appear in the Serial Monitor but the hardware does not respond, the issue is likely related to the hardware wiring or firmware configuration rather than the Bluetooth connection.

\section{Customizing Telemetry Labels in the Android Application}

The Android application automatically displays labels for telemetry
values such as numeric panels, indicators, and plots.

These labels are defined in the ESP32 firmware and sent to the
application when the Bluetooth connection is established.

This allows each project to customize the interface without modifying
the Android application.

\subsection{Where the Labels Are Defined}

Telemetry labels are defined inside the file:

\begin{center}
	\texttt{telemetry.cpp}
\end{center}

Inside the function responsible for sending the telemetry configuration,
you will find the following variables:

\begin{verbatim}
	const char* plotNames[] = { "Volts", "Amps", "RPMs" };
	const char* panelNames[] = { "Left", "Right" };
	const char* indicatorNames[] = { "Throttle", "Battery" };
\end{verbatim}

These arrays define the labels that appear in the Android application's
RC screen.

\subsection{Plot Labels}

The \texttt{plotNames} array defines the names of the real-time graphs.

Example:

\begin{verbatim}
	const char* plotNames[] = { "Volts", "Amps", "RPMs" };
\end{verbatim}

In the Android application this will create three plots labeled:

\begin{itemize}
	\item Volts
	\item Amps
	\item RPMs
\end{itemize}

You can rename them to match your sensors.

Example for a robot car:

\begin{verbatim}
	const char* plotNames[] = { "Speed", "Motor Temp", "Battery Current" };
\end{verbatim}

\subsection{Numeric Panel Labels}

The \texttt{panelNames} array defines the labels for the large numeric
displays.

Example:

\begin{verbatim}
	const char* panelNames[] = { "Left", "Right" };
\end{verbatim}

Typical uses include:

\begin{itemize}
	\item Motor speed
	\item Distance
	\item RPM
	\item Temperature
\end{itemize}

Example modification:

\begin{verbatim}
	const char* panelNames[] = { "Motor RPM", "Battery Voltage" };
\end{verbatim}

Note: It's recommended to use a maximum of 12 characters. 

\subsection{Indicator Labels}

The \texttt{indicatorNames} array defines the labels for the indicator
widgets.

Example:

\begin{verbatim}
	const char* indicatorNames[] = { "Throttle", "Battery" };
\end{verbatim}

Possible uses include:

\begin{itemize}
	\item Throttle position
	\item Battery level
	\item Temperature
	\item Signal strength
\end{itemize}

Example modification:

\begin{verbatim}
	const char* indicatorNames[] = { "Engine Load", "Battery Level" };
\end{verbatim}

\subsection{When the Labels Are Sent}

The telemetry labels are transmitted only once after the Android
application connects to the ESP32.

The firmware sends a telemetry configuration packet that includes:

\begin{itemize}
	\item plot labels
	\item numeric panel labels
	\item indicator labels
\end{itemize}

After receiving this packet, the Android application automatically
updates the RC screen interface.

\subsection{Important Notes}

\begin{itemize}
	\item The number of labels must match the number of telemetry variables.
	\item After modifying the labels you must upload the firmware again.
	\item The Android application will update the labels automatically after reconnecting.
\end{itemize}


